\section{Finite strategic games: complete proofs}
\label{sec:annex-p17-finite-strategic-games}

\begin{annexitem}{annex:p17.mixed.payoffs}{thm:p17-mixed-payoffs}{Independence, expectation and permissible scales}
\begin{proof}
Every product mass is nonnegative, and distributivity gives
\[
       \sum_{a\in\mathcal A}\prod_jx_j(a_j)
                     =\prod_j\sum_{a_j\in A_j}x_j(a_j)=1.
\]
Thus \(P_x\) is a lottery on the finite outcome set.
Finite summation gives exactly the expected utility of
\cref{def:p17-expected-utility}. Each probability coordinate
occurs to degree at most one in each summand, so the expectation
is a continuous polynomial and is affine in each mixture
separately. Grouping terms by \(a_i\) proves the displayed
own-mixture identity. The same expansion works for any player
coordinate, whether or not it belongs to the player receiving
the payoff.

For uniqueness, express every mixture as
\(\sum_{a_i}x_i(a_i)\delta_{a_i}\) by the proved lottery
decomposition. Expand a separately affine extension in
coordinate 1, then coordinate 2, and so on. After finitely
many expansions the result is necessarily
\(\sum_a(\prod_jx_j(a_j))u_i(a)\). This also proves
uniqueness when a player has just one action.
When \(n=1\), the product over opponents is the empty
product 1, giving the ordinary one-person lottery expectation.

For \(\widetilde u_i=\alpha_i u_i+\beta_i\), normalization gives
\(\widetilde U_i(x)=\alpha_i U_i(x)+\beta_i\).
The difference between any two unilateral choices is multiplied
by the positive number \(\alpha_i\), preserving its sign
and every tie. This proves best-response and equilibrium
invariance. The preference justification for precisely these
scales is \cref{thm:p17-vnm-affine-uniqueness}.
\end{proof}
\end{annexitem}

\begin{annexitem}{annex:p17.best.response.geometry}{thm:p17-best-response-geometry}{Every best-response and support condition}
\begin{proof}
Fix the opponents. Let \(w_a=U_i(a,x_{-i})\) and
\(v=\max_a w_a\), which exists because \(A_i\) is finite
and nonempty. For a mixture \(y_i\),
\[
              v-U_i(y_i,x_{-i})=\sum_a y_i(a)(v-w_a).
\]
All summands are nonnegative. The sum is zero exactly when
each action of positive weight has \(w_a=v\).
Thus the maximizers are precisely all lotteries on the
nonempty maximizing action set. This is a nonempty closed
convex simplex face. In particular checking pure deviations
checks all mixed deviations, and the three stated equilibrium
criteria are equivalent.

The space \(X\) is a nonempty closed bounded subset of
a finite-dimensional Euclidean space, given by nonnegative
coordinates and one unit-sum equation per player. It is
convex and compact. Each \(\mathcal B(x)\) is a finite
product of compact convex nonempty faces and has the
same properties.
To prove closed graph, suppose \(x^m\to x\),
\(y^m\to y\), and \(y^m\in\mathcal B(x^m)\).
For every player and pure action,
\[
                  U_i(y_i^m,x_{-i}^m)\geq U_i(a_i,x_{-i}^m).
\]
Continuity lets these inequalities pass to the limit.
The pure-deviation criterion makes \(y_i\) a best response
to \(x_{-i}\), proving \(y\in\mathcal B(x)\).
Sequential closedness is closedness in this finite-dimensional
metric space.

Finally \(\mathcal E\) is the intersection inside \(X\)
of the finitely many closed sets
\(\{x:U_i(x)-U_i(a_i,x_{-i})\geq0\}\).
These defining functions are polynomials. This proves
closedness, compactness, and the finite polynomial
description without presupposing that the set is nonempty.
\end{proof}
\end{annexitem}

\begin{annexitem}{annex:p17.sperner.foundation}{thm:p17-sperner-foundation}{The entire subdivision and parity construction}
\begin{proof}
A geometric \(m\)-simplex is the convex hull of \(m+1\)
affinely independent points; its faces are convex hulls
of subsets of those vertices. Write \(b_F\) for the
barycenter of a nonempty face \(F\). We first verify the
subdivision geometrically. In a simplex with vertices
\(v_1,\ldots,v_{m+1}\), write a point uniquely as
\(\sum_j\lambda_jv_j\), with \(\lambda_j\geq0\) and
\(\sum_j\lambda_j=1\).
Choose an ordering \(\sigma\) with
\(\lambda_{\sigma(1)}\geq\cdots\geq\lambda_{\sigma(m+1)}\),
and put \(\lambda_{\sigma(m+2)}=0\).
Let \(F_k\) contain the first \(k\) ordered vertices and set
\[
                  \mu_k=k(\lambda_{\sigma(k)}
                                      -\lambda_{\sigma(k+1)}).
\]
These numbers are nonnegative and telescope to
\(\sum_k\mu_k=\sum_j\lambda_j=1\). Moreover
\[
 \sum_{k=1}^{m+1}\mu_k b_{F_k}
   =\sum_{j=1}^{m+1}
         \left(\sum_{k=j}^{m+1}\frac{\mu_k}{k}\right)v_{\sigma(j)}
   =\sum_j\lambda_{\sigma(j)}v_{\sigma(j)}.
\]
Thus the simplexes associated with full chains of faces
cover the original simplex. Their barycenter vertices
are affinely independent: the inverse formulas
\(\lambda_{\sigma(j)}=\sum_{k\geq j}\mu_k/k\)
make their barycentric coordinates unique.

The faces with positive \(\mu_k\) are precisely the
strict upper-level sets of the coordinate list
\((\lambda_j)\). They do not depend on how equal
coordinates were ordered. Hence a point in two
chain simplexes uses, with positive weight, exactly
barycenters common to both chains. Their intersection
is therefore the convex hull of their common
barycenter vertices, a common face. On an original
boundary face all omitted coordinates are zero;
the formula uses only subfaces of that boundary
with positive coefficients. Subdivisions of adjacent
original simplexes consequently agree on their
common face. There are finitely many face chains.
This proves every triangulation and boundary
compatibility assertion.

Let \(\delta\) be the diameter of an original
\(m\)-simplex. For nested nonempty faces \(F\subsetneq G\),
let \(p\) and \(q\) be their numbers of vertices, and let
\(b_{\rm rem}\) be the average of the vertices of \(G\)
not in \(F\). Then
\[
       b_G=\frac{p}{q}b_F+\frac{q-p}{q}b_{\rm rem}.
\]
Both barycenters on the right lie in the original
simplex, whose diameter is \(\delta\), so
\[
             \|b_G-b_F\|
                 \leq\left(1-\frac{p}{q}\right)\delta
                 \leq\frac{m}{m+1}\delta.
\]
The diameter of a convex hull of finitely many points
is at most their largest pairwise distance: expand
the difference of two convex combinations as a
weighted sum of pairwise differences and use the
triangle inequality. Thus this bound holds for
every subdivided simplex. For all \(m\leq d\)
it is at most \(d/(d+1)\) times the preceding
mesh. Iteration proves the shrinking-mesh claim.
The zero-dimensional case is a single point.

We now prove the oddness assertion by induction
on \(d\). For \(d=0\), the unique vertex must have
label 0. For \(d\geq1\), call a \((d-1)\)-face
special if its \(d\) vertex labels are exactly
\(0,\ldots,d-1\).
Count incidences of special faces with top-dimensional
simplexes. An interior facet has exactly two
incident top-dimensional simplexes; a boundary
facet has exactly one. To justify this elementary
geometric fact, take a point in the facet's relative
interior and a sufficiently small neighborhood
avoiding every other face not containing it.
The facet's affine hyperplane has two local sides.
A top-dimensional simplex with that facet fills
one of these sides near the chosen point.
Disjoint simplex interiors allow at most one
on either side, and coverage requires both
sides for an interior point, or the sole available
side for a boundary point. This also proves
the assertion for \(d=1\), where the facets
are subdivision points of an interval.

A boundary special facet can lie only in the
outer face \(z_d=0\). Any other outer facet
forbids one of its required labels. A
\((d-1)\)-dimensional simplex cannot lie in
an intersection of two distinct outer facets,
which has dimension at most \(d-2\).
The induced triangulation of \(z_d=0\) satisfies
the \((d-1)\)-dimensional labeling condition.
By induction it has an odd number of special
facets. Interior incidences occur twice, so
the total incidence count is odd.

Inspect a top-dimensional simplex's labels.
If it has all \(d+1\) labels, exactly one
facet is special, the facet omitting label \(d\).
If it has all labels \(0,\ldots,d-1\) but
no label \(d\), exactly one label occurs twice,
and dropping either occurrence gives a special
facet: its contribution is two.
Every other label pattern contributes zero.
Consequently the incidence count modulo two
equals the number of fully labeled simplexes
modulo two. The former is odd, proving the
parity lemma and its nonemptiness conclusion.
\end{proof}
\end{annexitem}

\begin{annexitem}{annex:p17.brouwer.foundation}{thm:p17-brouwer-foundation}{The labeling limit and the convex-domain projection}
\begin{proof}
We use the elementary compactness of closed bounded
subsets of Euclidean space. Its sequential form
follows by repeatedly bisecting a bounding interval
in each coordinate, retaining a box containing
infinitely many sequence terms, and choosing a
diagonal subsequence. The box diameters tend to
zero, and completeness of the real numbers gives
a limit; a closed set contains that limit.
Continuous functions on such a set are bounded
and attain their infima and suprema: otherwise
an unbounded-value sequence, or an approximating
sequence for a finite extremum, contradicts
continuity along a convergent subsequence.

First let \(f:\Delta_d\to\Delta_d\) be continuous.
For \(d=0\) its unique point is fixed.
For \(d\geq1\), take the successive barycentric
triangulations. If a vertex in any triangulation
is fixed, the conclusion follows. Otherwise
label each vertex \(v\) by an index \(j\) with
\(v_j>f_j(v)\), choosing the least such index.
One exists because the coordinate sums of
\(v\) and \(f(v)\) are both 1 and they differ.
It has \(v_j>0\), so this is a Sperner labeling.
Choose a fully labeled simplex from each
triangulation, with vertex \(v_j^m\) bearing
label \(j\).

Compactness gives a subsequence of \(v_0^m\)
converging to some \(z\in\Delta_d\). Shrinking
mesh makes every \(v_j^m\) along the same
subsequence converge to \(z\).
Passing \((v_j^m)_j>f_j(v_j^m)\) to the
limit gives \(z_j\geq f_j(z)\) for all \(j\).
Their differences sum to zero, so all
differences are zero and \(f(z)=z\).

For the general convex domain, define the
nearest-point projection \(\pi_K(x)\) for
each \(x\in\mathbb R^d\). A minimizer of
\(\|x-z\|\), \(z\in K\), exists by compactness.
It is unique: if distinct \(p,q\) both
minimize the squared distance, their midpoint
is in \(K\) and
\[
 \left\|x-\frac{p+q}{2}\right\|^2
   =\frac{\|x-p\|^2+\|x-q\|^2}{2}-\frac{\|p-q\|^2}{4}
\]
is strictly smaller. For \(p=\pi_K(x)\) and
any \(z\in K\), minimality along
\(p+t(z-p)\), \(0\leq t\leq1\), gives
\[
                         (x-p)\cdot(z-p)\leq0.
\]
Indeed expansion of the squared distance,
subtraction of its value at \(t=0\), and
division by \(t>0\) give this inequality
as \(t\) decreases to zero.
For \(q=\pi_K(y)\), use \(z=q\) in the
inequality for \(x,p\) and \(z=p\) in
the one for \(y,q\). Subtraction yields
\[
                \|p-q\|^2\leq(x-y)\cdot(p-q)
                                   \leq\|x-y\|\|p-q\|.
\]
The last inequality is Cauchy--Schwarz,
obtained by minimizing the nonnegative
quadratic \(\|v-tw\|^2\).
Hence \(\pi_K\) is 1-Lipschitz and in
particular continuous. It fixes every point
of \(K\).

For \(d\geq1\), choose \(R>0\) with
\(K\subseteq[-R,R]^d\). The simplex
\[
          S=\{x:x_j\geq-R,\ \sum_{j=1}^d x_j\leq dR\}
\]
contains that box. Its vertices are
\((-R,\ldots,-R)\) and this vector plus
\(2dR e_j\), for \(1\leq j\leq d\);
it is an affine copy of \(\Delta_d\).
The continuous map \(f\circ\pi_K:S\to K\subseteq S\)
has a fixed point \(x\) by the simplex case.
Then \(x\in K\), so \(\pi_K(x)=x\) and
\(f(x)=x\). This covers lower-dimensional
and singleton \(K\) without an interior
assumption. In \(\mathbb R^0\) there is
only one point, so the conclusion is immediate.
\end{proof}
\end{annexitem}

\begin{annexitem}{annex:p17.kakutani.foundation}{thm:p17-kakutani-foundation}{Finite approximate selections and the full limiting argument}
\begin{proof}
We first extract the needed consequence of closed graph.
For a nonempty set \(A\), write
\(\operatorname{dist}(y,A)=\inf_{a\in A}\|y-a\|\).
For a fixed \(x\in K\) and \(\eta>0\), there is a
\(\delta>0\) such that
\[
 z\in K,\ \|z-x\|<\delta,\ y\in\mathcal C(z)
                    \quad\Longrightarrow\quad
                    \operatorname{dist}(y,\mathcal C(x))<\eta.
\]
Otherwise choose \(z_m\to x\) and
\(y_m\in\mathcal C(z_m)\) with these distances
at least \(\eta\). A subsequence of \(y_m\)
converges in compact \(K\) to \(y\).
Closed graph gives \(y\in\mathcal C(x)\),
while the distances are at most \(\|y_m-y\|\)
along that subsequence, a contradiction.
Each \(\mathcal C(x)\) is compact because
it is closed in \(K\).

For \(\epsilon>0\), choose a finite
\(\epsilon/2\)-net \(z_1,\ldots,z_m\) in \(K\).
Such a net follows by dividing a bounding
box into finitely many boxes of sufficiently
small diameter and taking a point of \(K\)
from each nonempty box. Choose
\(w_j\in\mathcal C(z_j)\), and put
\[
 \theta_j(x)=\max\{0,\epsilon-\|x-z_j\|\},\qquad
 f_\epsilon(x)=\frac{\sum_j\theta_j(x)w_j}
                         {\sum_j\theta_j(x)}.
\]
At least one \(\theta_j(x)\) is at least
\(\epsilon/2\), so the denominator is positive.
This is a continuous map \(K\to K\),
since its value is a convex combination
of points of \(K\). Brouwer therefore
gives \(x_\epsilon=f_\epsilon(x_\epsilon)\).

Take \(\epsilon=1,1/2,1/3,\ldots\) and
a convergent subsequence of the resulting
fixed points, with limit \(x\in K\).
If a weight at \(x_\epsilon\) is positive,
its center satisfies
\(\|z_j-x_\epsilon\|<\epsilon\).
Thus all centers contributing to this
fixed point tend uniformly to \(x\).
For the \(\delta\) given above and sufficiently
late terms, every contributing \(w_j\)
is within \(\eta\) of \(\mathcal C(x)\).
Choose a nearest \(v_j\in\mathcal C(x)\)
for each contributing \(w_j\), possible
by compactness. The same normalized
weighted average of the \(v_j\)'s belongs
to \(\mathcal C(x)\) by convexity, and
is within \(\eta\) of \(x_\epsilon\).
It follows that
\[
 \operatorname{dist}(x,\mathcal C(x))
                           \leq\|x-x_\epsilon\|+\eta.
\]
Let the subsequence tend to its limit
and then let \(\eta\) tend to zero.
Closedness of \(\mathcal C(x)\) gives
\(x\in\mathcal C(x)\), as required.
No continuous exact selection was presumed.
\end{proof}
\end{annexitem}

\begin{annexitem}{annex:p17.nash.1950}{thm:p17-nash-1950}{The complete original countering route}
\begin{proof}
The product \(X\) is nonempty, compact and convex.
The best-response lemma gives every value
\(\mathcal B(x)\) nonempty, compact and convex,
and gives its graph closed. The fully proved
correspondence fixed-point theorem therefore
gives \(x\in\mathcal B(x)\).
By definition this means that each component
is optimal against the fixed other components,
which is Nash equilibrium.
Conversely any equilibrium counters itself
by the same definition. This is precisely
the correspondence used in the 1950 note,
with every required fixed-point premise
and its underlying proof supplied.
\end{proof}
\end{annexitem}

\begin{annexitem}{annex:p17.nash.1951}{thm:p17-nash-1951}{The original positive-gain map in both directions}
\begin{proof}
Each \(g_{i,a}\) is continuous and nonnegative,
and the denominator \(1+G_i(x)\) is at least 1.
The new coordinates are nonnegative and sum to
\((1+G_i(x))/(1+G_i(x))=1\).
Thus \(T\) is a continuous self-map of \(X\).

If \(x\) is an equilibrium, every unilateral
pure gain is nonpositive, so every \(g_{i,a}\)
is zero and \(T(x)=x\).
Conversely suppose \(T(x)=x\), and fix a player.
Its positive support is nonempty and its
expected payoff is the weighted average
of the supported pure payoffs.
At least one supported action \(a\)
has payoff no greater than that average:
if all were greater, their weighted
average would be greater than itself.
For this action \(x_i(a)>0\) and
\(g_{i,a}(x)=0\).
The fixed-coordinate equation is
\[
                     x_i(a)G_i(x)=g_{i,a}(x)=0.
\]
It forces \(G_i(x)=0\), and hence every
nonnegative \(g_{i,b}(x)=0\).
Repeating for each player gives all pure
deviation inequalities, hence equilibrium.
This includes a one-action player directly.
Brouwer on \(X\) now supplies a fixed point
and proves existence by the 1951 construction.
The argument asserts no orbit-convergence
property of this continuous map.
\end{proof}
\end{annexitem}

\begin{annexitem}{annex:p17.nash.symmetry}{thm:p17-nash-symmetry}{All automorphisms and the invariant convex domain}
\begin{proof}
The profile assigning equal probability to
every action of each player is invariant.
An automorphism sends a player's action
set bijectively onto the target player's
equally sized set, so it preserves these
uniform probabilities.
Each automorphism acts on the ambient
profile coordinates by a linear permutation.
Its fixed set is a closed linear subspace;
intersecting these fixed sets with \(X\)
gives a closed convex subset \(K\) of
compact \(X\), containing the uniform
profile.

Reindexing the finite payoff sum gives
\[
              U_{\sigma(i)}(h(x))=U_i(x).
\]
The same reindexing with a unilateral pure
replacement gives
\[
       g_{\sigma(i),h_i(a)}(h(x))=g_{i,a}(x).
\]
The gain sums and probability numerators
are permuted in the same way, so
\(T(h(x))=h(T(x))\).
Consequently \(T\) preserves the common
fixed set \(K\). The general convex-domain
Brouwer theorem applies even if \(K\)
has lower dimension and supplies a
fixed point there. The 1951 map theorem
makes it an equilibrium, and membership
in \(K\) makes it invariant under every
automorphism. This proves the original
symmetry conclusion with its exact scope.
\end{proof}
\end{annexitem}

\begin{annexitem}{annex:p17.nash.subsolutions}{thm:p17-nash-subsolutions}{Maximal interchangeable products and all limiting cases}
\begin{proof}
The existence theorem gives \(\mathcal E\ne\varnothing\).
Fix an equilibrium \(x\). Its singleton is interchangeable.
Order by inclusion the interchangeable subsets containing it.
The union of a nonempty chain is interchangeable: any two
profiles in the union occur together in one chain member,
where every required coordinate replacement already belongs.
The union remains a subset of \(\mathcal E\).
The empty chain has the singleton as an upper bound.
Zorn's lemma therefore gives a maximal member, which is
a subsolution in the full family.

For any nonempty interchangeable \(C\), choose one profile
from \(C\) to start. Given \(z_i\in C_i\) for every \(i\),
some member of \(C\) has component \(z_i\).
Successively replace coordinates \(1,\ldots,n\) of the
starting profile by these components.
Interchangeability keeps every intermediate profile in
\(C\), and the final profile is \(z\).
Thus \(\prod_iC_i\subseteq C\); the reverse inclusion
is the definition of the factors. This proves the product
claim before using maximality.

Let now \(C\) be a subsolution, and fix \(i\).
Take \(s_i,t_i\in C_i\), \(0\leq\lambda\leq1\), and
put \(z_i=\lambda s_i+(1-\lambda)t_i\).
For any \(x_{-i}\in\prod_{j\ne i}C_j\), both completed
profiles with \(s_i\) and \(t_i\) are equilibria.
For every player \(j\) and pure deviation \(a_j\),
the difference between its current payoff and its
deviation payoff is affine in the component being
varied. When \(j=i\), the deviation term is constant
in that component; when \(j\ne i\), both terms
are affine in it. The same difference at \(z_i\)
is therefore a convex combination of two nonnegative
differences and is nonnegative.
So every profile in
\((C_i\cup\{z_i\})\times\prod_{j\ne i}C_j\)
is an equilibrium. That product is interchangeable;
maximality forces \(z_i\in C_i\). This proves full
convexity, including mixture endpoints.

If \(s_i^m\in C_i\) converges to \(z_i\), fix again
any other profile in the product of the other
factors. The completed profiles are equilibria
and converge to the completed profile with \(z_i\).
Closedness of \(\mathcal E\) puts that limit in
\(\mathcal E\). Adding \(z_i\) to the factor thus
again gives an interchangeable equilibrium product,
and maximality puts \(z_i\) in \(C_i\).
Each factor is closed in its compact simplex,
hence compact. It is nonempty because \(C\) is.
Their finite product is compact and equals \(C\).

If the entire \(\mathcal E\) is interchangeable,
it contains every interchangeable subset and
is the sole maximal one. Conversely, if there
is a unique subsolution, every equilibrium
belongs to it by the first paragraph. It is
therefore the entire \(\mathcal E\), proving
solvability. Continuity on the nonempty compact
sets \(\mathcal E\) and \(C\) gives attained
upper and lower payoff values. A single
equilibrium gives identical maximum and minimum.
The converse is unnecessary: constant payoff
games can have many equilibria with the
same payoff vector.
\end{proof}
\end{annexitem}

\begin{annexitem}{annex:p17.nash.polyhedral}{thm:p17-nash-polyhedral}{The source slice construction and the full vertex argument}
\begin{proof}
In a solvable game, the product theorem gives
\(\mathcal E=\prod_i C_i\). For fixed \(t\in\mathcal E\),
\[
                     z_j\in C_j
                    \quad\Longleftrightarrow\quad
                     t[j\leftarrow z_j]\in\mathcal E.
\]
The pure-deviation inequalities for the profile
on the right are exactly those in the statement.
The current payoff is affine in \(z_j\).
If \(i=j\), the deviated profile has replaced
that coordinate completely and its payoff is
constant in \(z_j\). If \(i\ne j\), the
deviation fixes another coordinate and leaves
an affine dependence on \(z_j\).
Thus every constraint is affine in \(z_j\).
Together with nonnegativity and the unit-sum
equation, finitely many such constraints
define the factor as a bounded polyhedron.

We supply the finite-vertex conclusion rather
than invoking it without proof. Let \(P\) be
any nonempty bounded intersection of finitely
many closed affine half-spaces in its affine
hull \(V\), of dimension \(r\).
Call an inequality active at a point when
equality holds there, and restrict its
linear normal to the direction space of \(V\).
A point is a vertex exactly when its active
normals span that direction space.
For if they do not span it, choose a nonzero
direction annihilating all of them.
The finitely many inactive inequalities have
positive slack, so sufficiently small motion
in either direction remains feasible.
The point is then the midpoint of distinct
feasible points and is not a vertex.
Conversely, if the normals span and a point
is the midpoint of two feasible points,
every active inequality must be equality
at both endpoints. Their difference
annihilates all active normals and is zero.
Thus the point is a vertex.

Each vertex is uniquely determined within
\(V\) by \(r\) independent active equations.
There are only finitely many choices from
the finite constraint list, so there are
only finitely many vertices.
Every point is their convex combination,
as follows. At a nonvertex \(x\), choose
an annihilating direction as above.
The maximal feasible interval on its line
has endpoints \(x+\alpha h,x+\beta h\),
where \(\alpha<0<\beta\). The endpoints
exist because \(P\) is closed and bounded.
At each endpoint an additional inequality
with nonzero value on \(h\) becomes active,
so the active-normal rank strictly increases.
Moreover
\[
 x=\frac{\beta}{\beta-\alpha}(x+\alpha h)
                      +\frac{-\alpha}{\beta-\alpha}(x+\beta h).
\]
Repeat at any nonvertex endpoint. The
rank can increase at most \(r\) times,
so the resulting finite binary decomposition
terminates at vertices. It expresses \(x\)
as their convex combination.
If \(r=0\), \(P\) is a singleton and
the assertion is immediate.
Applying this argument to each \(C_j\)
proves the source's stated vertex corollary,
including lower-dimensional factors.
\end{proof}
\end{annexitem}

\begin{annexitem}{annex:p17.finite.minimax}{thm:p17-finite-minimax}{Saddle inequalities, attained values and exact factor sets}
\begin{proof}
Let \(f(x,y)=x^{\mathsf T}My\).
For any fixed \(x\), its minimum over \(y\)
is the minimum of its finitely many pure-column
values. For fixed \(y\), its maximum over
\(x\) is the maximum of finitely many
pure-row values. These finite minima and
maxima of continuous affine functions are
continuous. Compactness of the two simplexes
therefore makes the outer extrema attained.
Always
\[
           \max_x\min_y f(x,y)\leq\min_y\max_x f(x,y),
\]
because \(\min_{y'}f(x,y')\leq f(x,y)\leq
\max_{x'}f(x',y)\) for every \(x,y\).

The finite-game theorem gives an equilibrium
\((x^*,y^*)\). Put \(v=f(x^*,y^*)\).
Player 1 maximizes \(f\) and player 2
maximizes \(-f\), so their deviation
conditions are precisely
\[
                    f(x,y^*)\leq v\leq f(x^*,y)
                                      \quad\hbox{for all }x,y.
\]
They show the left outer optimum is
at least \(v\), the right one at most
\(v\), and the preceding weak inequality
forces equality of both to \(v\).

A strategy in \(X^*\) guarantees at
least \(v\) against every pure column,
hence against every mixed column.
A strategy in \(Y^*\) holds the row
payoff at most \(v\) against every row
mixture. For any pair in \(X^*\times Y^*\)
the payoff is thus exactly \(v\);
the same inequalities show that neither
player can gain by deviating. Conversely
any equilibrium supplies the displayed
saddle inequalities and therefore has
the common minimax value \(v\), with its
components in \(X^*,Y^*\).
This proves the exact product description
and equality of all equilibrium payoffs.

If instead the attained minimax equality
is given first, choose maximizing and
minimizing optimal strategies for its
two sides. Their guarantees give the
same saddle inequalities and hence
an equilibrium. This proves the reverse
implication and the claimed equivalence.
The equilibrium product is interchangeable,
so the game is solvable.

In any strongly solvable finite game,
start at an equilibrium and replace
one player's component by a pure best
response, which exists. Its own payoff
is unchanged, so strong solvability
keeps the profile an equilibrium.
Repeat for every player, obtaining a
pure equilibrium. For the two-player
zero-sum game this is a pure saddle.
The converse fails for
\[
                            M=\begin{pmatrix}1&1\\1&0\end{pmatrix}.
\]
The pair (row 1, column 1) is a saddle,
as is (row 1, column 2).
Against column 1, row 2 also gives
payoff 1 and is a best response.
But (row 2, column 1) is not an
equilibrium: the column player can
switch to column 2 and lower the
row payoff to zero. Thus the
game has a pure saddle and is
not strongly solvable.
\end{proof}
\end{annexitem}

\begin{annexitem}{annex:p17.nash.six.games}{thm:p17-nash-six-games}{All supports and payoffs of the six original examples}
\begin{proof}
Let \(D_1(q)\) be row \(a\)'s payoff
minus row \(b\)'s payoff, and \(D_2(p)\)
column \(\alpha\)'s payoff minus column
\(\beta\)'s payoff. Direct evaluation
of every source matrix gives
\[
\begin{array}{c|cc}
\text{game}&D_1(q)&D_2(p)\\ \hline
1&17q-7&9-16p\\
2&-9&-9\\
3&22q-11&22p-11\\
4&0&0\\
5&8q-3&8p-2\\
6&q&p.
\end{array}
\]
The complete best-response rule is:
\(p=1\) if \(D_1(q)>0\), \(p=0\)
if \(D_1(q)<0\), and any \(p\)
is possible if \(D_1(q)=0\).
The identical rule with \(q,D_2\)
applies to the column player.
This includes every pure boundary and
every mixed support.

For game 1, if \(q>7/17\), then
\(p=1\), forcing \(D_2(p)<0\)
and \(q=0\), a contradiction.
If \(q<7/17\), then \(p=0\),
forcing \(q=1\), also a contradiction.
Thus \(q=7/17\). Since this is
strictly between zero and one,
\(D_2(p)=0\) and \(p=9/16\).
At this profile the row's two
pure payoffs are \(9q-4\) and
\(3-8q\), both \(-5/17\).
The column's are \(5-8p\) and
\(8p-4\), both \(1/2\).
This proves uniqueness and absence
of any pure equilibrium.

For game 2, both differences are
strictly negative everywhere, forcing
\(p=q=0\), with payoffs \((-1,-1)\).
The sole best response against that
profile is the same pure action for
each player, so the singleton solution
is strong.

For game 3, \(q>1/2\) forces
\(p=1\), then \(q=1\); \(q<1/2\)
forces \(p=0\), then \(q=0\).
At \(q=1/2\), an interior \(q\)
requires \(p=1/2\).
Thus exactly the two matching corners
and the midpoint occur. The corners
give \((1,1)\); averaging the four
equally likely entries at the midpoint
gives \((-9/2,-9/2)\).
For completeness, against pure
opponents either player's two payoff
expressions are \(11t-10\) and
\(1-11t\). Their minimum is maximized,
and their maximum minimized, at
\(t=1/2\), with value \(-9/2\).
This proves the source's additional
maximin/minimax observation about that
strategy; the actual players still
have identical, rather than opposite,
payoffs.

For game 4, each player is indifferent
between its own actions at every
opponent mixture. The payoffs are
\((q,p)\), so every profile is an
equilibrium and the whole product
is a strong solution. The upper
and lower equilibrium values are
1 and 0 for both players.

For game 5, the same threshold argument
as in game 3 uses \(q=3/8\) and
\(p=1/4\). A larger \(q\) forces
the matching first-action corner,
and a smaller \(q\) the matching
second-action corner; equality
forces the stated interior \(p\).
These are the only equilibria.
Their pure payoffs are the corresponding
matrix entries. At the interior
profile the row payoffs are
\(2q-1=2-6q=-1/4\), and the
column payoffs are
\(3p-1=1-5p=-1/4\).

For game 6, if \(q>0\), then
\(p=1\) and hence \(q=1\).
If \(q=0\) and \(p>0\),
the column rule would instead
force \(q=1\), a contradiction.
Thus the only other profile
is \(p=q=0\).
At \((\varepsilon,\varepsilon)\)
the current payoff is \(\varepsilon^2\);
certainly choosing the first action
gives \(\varepsilon\), a gain
\(\varepsilon(1-\varepsilon)\).

Any singleton equilibrium set is
interchangeable, proving solvability
of games 1 and 2, and the full
product proves it for game 4.
In each of games 3, 5 and 6,
interchanging one component between
the two matching pure equilibria
produces an unmatched corner which
is not in the complete equilibrium
list. Hence they are not solvable.
Game 1 is not strongly solvable:
against \(q=7/17\), every row
mixture is a best response, but
changing \(p\) from \(9/16\) while
keeping that interior \(q\) destroys
equilibrium. This finishes all
claimed classifications.
\end{proof}
\end{annexitem}

\begin{annexitem}{annex:p17.strategic.dominance}{thm:p17-strategic-dominance}{The finite comparison and the correctly directed perturbation}
\begin{proof}
The difference between the two strategies'
payoffs against a mixed opponent profile
is the weighted average, under the
product opponent law, of their differences
against pure opponent profiles.
If all the latter differences are strictly
positive, their finite minimum is positive,
and so is every such average.
The converse holds because pure profiles
are among the mixed profiles.
At an equilibrium a strictly dominating
replacement would be a profitable
unilateral deviation, proving exclusion.

Strict domination implies \(x_i\ne y_i\).
Since their coordinate sums agree,
there is at least one coordinate
with \(x_i(a)>y_i(a)\).
The hypothesis gives strictly positive
\(z_i(a)\) at every such coordinate,
so the displayed finite minimum is
strictly positive. For a positive
\(\rho\) no greater than that minimum,
every decreasing coordinate of
\(z_i+\rho(y_i-x_i)\) remains
nonnegative; every other coordinate
is nondecreasing. The sum remains
1. The resulting vector is a
lottery, and affinity gives
\[
 U_i(z_i',w_{-i})-U_i(z_i,w_{-i})
        =\rho[U_i(y_i,w_{-i})-U_i(x_i,w_{-i})]>0
\]
for every opponent profile.
Thus \(z_i'\), rather than \(z_i\),
is the dominating vector.

Let \(z_i\) be undominated and let
\(x_i\) have support contained in
that of \(z_i\).
If \(x_i\) were dominated by \(y_i\),
every coordinate with \(x_i(a)>y_i(a)\)
would have \(x_i(a)>0\), hence
\(z_i(a)>0\). The perturbation just
proved would then dominate \(z_i\),
a contradiction. Therefore the
entire face spanned by the support
of an undominated strategy is
undominated. Taking the union of
these faces gives precisely the
undominated set.
\end{proof}
\end{annexitem}

\begin{annexitem}{annex:p17.product.and.correlation}{thm:p17-product-and-correlation}{All recommendation and product-law cases}
\begin{proof}
A randomized recommendation-dependent deviation
is given by probabilities \(d_i(b_i\mid a_i)\)
on replacement actions, for each recommendation.
The expected payoff from obedience minus
that from the deviation equals
\[
 \sum_{a_i,b_i}d_i(b_i\mid a_i)
 \sum_{a_{-i}}\mu(a_i,a_{-i})
       [u_i(a_i,a_{-i})-u_i(b_i,a_{-i})].
\]
Every coefficient in the outer sum
is nonnegative. Thus the stated
inequalities imply nonnegative
loss from any such deviation.
Conversely, switch to one fixed
\(b_i\) on one specified recommendation
\(a_i\), and obey on all other
recommendations. Its loss is exactly
that one inequality, proving necessity.
The calculation includes zero-probability
recommendations, whose inner sums are zero.

For \(\mu=P_x\), the inner sum factors as
\[
 x_i(a_i)
       [U_i(a_i,x_{-i})-U_i(b_i,x_{-i})].
\]
If \(x\) is an equilibrium, every
positive-weight action is optimal,
so this expression is nonnegative.
If the product obeys every inequality,
each action of positive weight
has at least the payoff of every
other action and is optimal.
The support criterion then gives
Nash equilibrium. This proves both
directions without dividing by a
zero probability.

For the diagonal law in game 3,
each recommendation reveals the
matching opposing action. Obedience
gives 1, and switching gives \(-10\);
the relevant nontrivial inequality
is \(11/2\geq0\) for each recommended
action and player. Hence the law
is obedient and gives payoff
\((1,1)\). Its two marginals are
uniform. Their product has each
of the four outcomes at probability
\(1/4\), whereas the recommendation
law has zero off-diagonal mass
and diagonal mass \(1/2\).
Thus it is not a product, and
the marginal product has the
already calculated payoff
\((-9/2,-9/2)\).
\end{proof}
\end{annexitem}

\begin{annexitem}{annex:p17.strategic.welfare.boundaries}{thm:p17-strategic-welfare-boundaries}{The stated welfare and solution-object separations}
\begin{proof}
Game 2 has the unique equilibrium
\((b,\beta)\), with payoff \((-1,-1)\).
The profile \((a,\alpha)\) gives
\((1,1)\), a strict improvement
for both players. It is nevertheless
not an equilibrium because each
player can gain by deviating.
This proves the failure of Pareto
efficiency at equilibrium.

A positive affine transformation of
one player's payoff preserves the
sign of every payoff difference
for that player. Doing this
independently for all players
therefore preserves Pareto comparisons
and, by the mixed-payoff lemma,
the equilibrium set.
In game 5 the two pure equilibrium
payoff vectors are \((1,2)\)
and \((2,1)\). Doubling player 1's
scale gives sums \(2+2=4\)
and \(4+1=5\). Doubling player 2's
scale instead gives \(1+4=5\)
and \(2+2=4\). Thus their aggregate
ranking reverses under scale changes
which leave every player's lottery
preference and the equilibrium
set unchanged.

The Nash definition quantifies over
one changed component of a strategy
profile in \(X\).
The NM definition in
\cref{def:p18-nm-solutions} instead
quantifies over a set of imputation
vectors, uses the coalition feasibility
inequality from
\cref{def:p18-cooperative-game},
and requires both internal stability
and domination of every omitted
imputation. Neither its objects nor
its admissible deviations are those
of the Nash definition.
The supplied Lucas counterexample
therefore does not conflict with
finite normal-form existence.
Likewise the source's interchangeability
definition concerns a property of
the entire equilibrium set; the
computed games 3, 5 and 6 have
equilibria but fail that property.
These explicit definitions and
examples prove the asserted
distinctions.
\end{proof}
\end{annexitem}

\begin{annexitem}{annex:p17.finite.nash.approximation}{thm:p17-finite-nash-approximation}{The rounding error and finite search proof}
\begin{proof}
If \(M=0\), all payoffs are zero
and every profile is an equilibrium.
If \(L=0\), every action set is
a singleton and the unique profile
is an equilibrium. This proves
the initial case \(ML=0\).

Otherwise let \(x^*\) be an equilibrium,
whose existence has been proved.
For a player with \(m\) actions,
round the first \(m-1\) probabilities
down to multiples of \(1/k\)
and give the final action the
remaining probability. The result
\(y_i\) is nonnegative and sums to 1.
The total downward error of the
first \(m-1\) coordinates is at
most \((m-1)/k\); the last
coordinate has exactly the opposite
error. Thus
\[
                  \|y_i-x_i^*\|_1\leq\frac{2(m-1)}{k},
 \qquad
                  \sum_i\|y_i-x_i^*\|_1\leq\frac{2L}{k}.
\]
With all other components fixed,
changing component \(j\) changes
any player's expected payoff by
at most \(M\|x_j-y_j\|_1\):
expand in that coordinate and
bound each conditional payoff
in absolute value by \(M\).
Changing components one at a time
therefore gives
\[
                    |U_i(x)-U_i(y)|
                         \leq M\sum_j\|x_j-y_j\|_1.
\]
The same bound applies to a pure
deviation payoff, which is independent
of the deviator's own mixture.
For every \(i,a_i\), its gain at
the rounded profile is consequently
at most its gain at \(x^*\)
plus \(2M\sum_j\|y_j-x_j^*\|_1\).
The former is nonpositive, so
\[
                U_i(a_i,y_{-i})-U_i(y)
                      \leq\frac{4ML}{k}\leq\varepsilon.
\]
Hence a satisfactory grid profile
exists.

For \(m\) actions the grid vectors
correspond to nonnegative integer
solutions of \(c_1+\cdots+c_m=k\).
Placing \(m-1\) separators among
\(k\) identical units gives
\({k+m-1\choose m-1}\) solutions.
The players' grid choices combine
independently, giving the stated
finite product count.
For rational data every expected
payoff and comparison with
\(\varepsilon\) is rational and
can be evaluated exactly.
Exhaustive enumeration must reach
one of the satisfactory profiles
and terminate. This proves the
claimed accuracy and finite bound,
without claiming a small running
time or an exact irrational
equilibrium representation.
\end{proof}
\end{annexitem}
